\documentclass[10pt,a4paper]{amsart}
\usepackage[margin=19mm]{geometry}
\usepackage{amsmath,amssymb,mathtools}
\usepackage[scr=rsfs,cal=euler]{mathalfa}
\usepackage{xfrac}
\usepackage[unicode,hidelinks,bookmarks=false]{hyperref}
\newcommand{\ie}{i.e.\ }
\newcommand{\cf}{cf.\ }
\newcommand{\resp}{respectively\ }
\newcommand{\Subsection}{\subsection}
\providecommand{\nobreakdash}{\nobreak\hbox{-}}
\setlength{\emergencystretch}{2em}
\multlinegap0pt
\usepackage{mathtools}
\usepackage{amssymb}
\usepackage{xcolor}
\usepackage[shortlabels]{enumitem}
\usepackage{algorithm}\usepackage{algpseudocode}
\usepackage{tikz}
\newtheorem{lemma}{Lemma}
\newcommand{\cC}{\mathcal{C} }
\title{A linear upper bound for zero-sum Ramsey numbers of bounded degree graphs}
\author{Jasmin Katz, Xiaopan Lian, Alexandru Malekshahian, Andrey Shapiro}
\date{Source: arXiv:2512.17790v2 (CC BY 4.0)}
\begin{document}
\maketitle

\noindent\emph{Source and licence.} A linear upper bound for zero-sum Ramsey numbers of bounded degree graphs. Version arXiv:2512.17790v2 (CC BY 4.0), distributed by arXiv under the Creative Commons Attribution 4.0 International licence, http://creativecommons.org/licenses/by/4.0/, as recorded in the arXiv licence metadata for this exact version and shown on the abstract page https://arxiv.org/abs/2512.17790v2. Full licence notice: \texttt{licenses/arxiv-2512.17790-CC-BY-4.0.txt}.\par\smallskip

\noindent\textit{Editorial scope.} Counted unit: the article's lemma core\_phase\_arg (source0.tex lines 1186--1196). The article's complete original proof of it is delivered with it (source0.tex lines 1198--1235, one \texttt{proof} environment of 38 lines). No mathematics is written, repaired or re-derived: the statement and the proof are the article's own lines, in the article's own notation, and no retained line is substituted or re-flowed.  No further source-local context is needed: every cross-reference of every retained line resolves inside the two delivered spans.  Nothing was omitted from any retained span, so the package carries no omission record.  No retained line contains a citation, so the wrapper carries no bibliography and no bibliography entry.  No retained line contains a figure environment, an image-inclusion command or drawing code, so no image is delivered and the package declares no asset dependency.  Editorial formatting only, in addition: an AMS \texttt{amsart} preamble that restates the article's own declarations and macros used by the retained lines, namely the environments \texttt{lemma} on separate counters, together with the standard packages those lines need.  The retained environments print in this document as Lemma 1 in order of appearance, whereas the article prints the same items with the numbering of its own full document.  The complete native source of the article as distributed in the arXiv e-print for this exact version is bundled unchanged as \texttt{sources/191346/source0.tex}.
\begin{lemma}\label{core_phase_arg}
    Suppose that $d, d'\leq\Delta$ and $m\leq \min(d, d')-1$ are positive integers such that $\gcd(d, d', n)=\kappa$ and that   
    $X\subset R_0$ is some pool of vertices. Let $c=c_1/H':E(R_0)\to \Gamma/H'$ (for some $H'\leqslant \Gamma$) be an edge coloring and $\cC:R_0\to \Gamma_0$ a vertex coloring that is constant on $X$. 
    
    If there are no $(d,\lambda)$--, $(d',\lambda)$-- or $(d,d',m,\lambda)$--gadgets under the colorings $c$ and $\cC$, then there exist sets $I\subset X$, $S=\{s_1,\ldots,s_{\lambda-1}\}\subset \Gamma$, $T\subset \Gamma$, and a vertex coloring $\cC':I\to T$ such that: 
    \begin{enumerate}
        \item $|I|\geq |X|/(\lambda-1)^{d+d'} - d-d'$;
        \item for all $t\in T$, we have $\kappa t \in H'$; and
        \item for all $x,y\in I$ distinct, there exists some $s\in S$ such that $c(xy)=s+\cC'(x)+\cC'(y)+H'$.
    \end{enumerate}  
\end{lemma}

\begin{proof}
First, if $|X|\leq d+d'$ then setting $I=T=\emptyset$ the lemma is vacuously true. Otherwise, fix sets $U=\{u_1,\ldots,u_{d-1}\}$ and $V=\{v_1,\ldots,v_{d'-1}\}$ in $X$ so that $M\coloneqq U\cap V =\{u_1,\ldots,u_{m}\}=\{v_1,\ldots,v_{m}\}$. Since $X$ does not contain a $(d', \lambda)$--gadget, it follows that for each $1\leq i<d'$ we have
\[|\{c(Uw'v_i):w'\in X\setminus (U\cup V)\}|\leq \lambda-1.\]
Denote by $g_0^i$ the most popular color in the set above, so that there exist at least $(|X|-(d+d'-m-2))/(\lambda-1)$ vertices $w'\in X\setminus (U\cup V)$ such that $c(Uw'v_i)=g_0^i$.

Similarly, note that for each $1\leq j<d$
\[|\{c(u_jw'V):w'\in X\setminus (U\cup V)\}|\leq \lambda-1\]
holds and analogously define $q_0^j$ to be the most popular color in the above set.

Then there exists a set of vertices $I_1\subset X$ with $|I_1|\geq (|X|-(d+d'-m-2))/(\lambda-1)^{d+d'-2}$ such that for all $w'\in I_1$, $i \in [d'-1]$ and $j \in [d-1]$, we have $c(Uw'v_i)=g_0^i$ and $c(u_jw'V)=q_0^j$. 

Further, let us fix a vertex $w\in I_1$ and observe that $c(Uw'w)$ and $c(ww'V)$ both take at most $\lambda-1$ colors across all $w' \in I_1\setminus \{w\}$, and so we can again restrict to the most popular colors (denoted by $g_1$ and $q_1$, respectively) to obtain a set $I\subset I_1\setminus\{w\}$ of size
\[|I|\geq \left(|X|-(d+d'-m-2)\right)/(\lambda-1)^{d+d'}-1\geq \left(|X|/(\lambda-1)^{d+d'}\right) -d-d'\] 
such that for all $w'\in I$, $c(Uw'w)=g_1$ and $c(ww'V)=q_1$. Further, let us define $$g_0\coloneqq \sum_{i< d'}g_0^i = c(w'V) + (d'-1)\cdot c(Uw')$$ and $$q_0\coloneqq \sum_{j< d} q_0^j = c(w'U) + (d-1)\cdot c(w'V).$$

Then observe that for all $w'\in I$, we have
$$
\begin{aligned}
c(ww')&=c(ww') +c(w'U) - c(w'U) - (d-1)\cdot c(w'V)+(d-1)\cdot c(w'V)\\
&=g_1-q_0+(d-1)\cdot c(w'V)\\
&=g_1-q_0+(d-1)\cdot(q_1-c(ww'))\\
&= g_1-q_0+(d-1)\cdot q_1-(d-1)\cdot c(ww')
\end{aligned}$$
and thus $$d\cdot c(ww')=g_1-q_0+(d-1)\cdot q_1,$$
which is a constant in $\Gamma/H'$ independent of the choice of $w'$. Similarly, we also get that $$d'\cdot c(ww') =q_1-g_0+(d'-1)\cdot g_1$$
which is again independent of $w'$.

Then, for distinct $w',w''\in I$ we have that $d\cdot (c(ww')-c(ww''))=d'\cdot (c(ww')-c(ww''))=H'$ (recalling that $c=c_1/H'$ takes values in $\Gamma/H'$). Now, since $\gcd(d,d',n)=\kappa$ it follows by B\'ezout's identity that \[\kappa\cdot c(ww')=\kappa\cdot c(ww'')\]
for all $w', w''\in I$.

Next, fixing some $v_0\in I$, we define for all $x\in I$ the new vertex coloring $\cC':I\to\Gamma$ by letting $\cC'(x)$ be an arbitrary representative of the coset $c(xw)-c(wv_0)\in \Gamma/H'$, and note that $\kappa \cdot \cC'(x)\in H'$ for all $x\in I$. We then define $T\coloneqq \{\cC'(x):x\in I\}\subset \Gamma$ and note it satisfies the required properties.

We now only have left to define $S$ and ensure that it and $\cC'$ satisfy the third condition in the lemma. We observe that $c(UxyV)=c(xy)+q_1+g_1-c(xw)-c(yw)$, for all $x\neq y\in I$ and so
\[c(xy)=c(UxyV) -q_1-g_1+\cC'(x)+\cC'(y)+2c(v_0w)+H'.\] 
Now, by our assumption that a $(d,d',m,\lambda)$--gadget does not exist, we see that $c(UxyV)$ can take at most $\lambda-1$ colors as $x, y$ range in $I$; call this set of attained colors $S'$ and extend it so it has cardinality exactly $\lambda-1$, i.e. we have $S'=\{s'_1, \dots, s'_{\lambda-1}\}\subset \Gamma/H'$.

Letting $\hat{s}= 2c(v_0w)-q_1-g_1\in \Gamma/H'$ constant, we then have that $c(xy)=\hat{s}+c(UxyV)+\cC'(x)+\cC'(y)$. Then for each $i\in [\lambda-1]$ take some arbitrary coset representative $s_i$ of $s_i' + \hat{s}$ (we recall that $s_i', \hat{s}\in \Gamma/H'$) and define $S=\{s_1,\ldots,s_{\lambda-1}\}\subset \Gamma$ and conclude that for all $x\neq y\in I$ there is some $s_i\in S$ with $s_i+H'=\hat{s}+c(UxyV)$ and so  $c(xy)=s_i+\cC'(x)+\cC'(y)+H'$.
\end{proof}

\end{document}
